\section{Conservative Vector Fields}
\subsection{Potential Functions and the Fundamental Theorem}

\begin{definition}{Conservative Field and Potential}{conservative_field}
A continuous vector field $\vv{F}$ on an open domain $D$ is \textbf{conservative} if there is a differentiable scalar function $f$ on $D$ such that
\[
    \vv{F} = \nabla f.
\]
The function $f$ is a \textbf{potential function}. Adding a constant to $f$ gives another potential for the same field.
\end{definition}

\begin{theorem}{Fundamental Theorem for Line Integrals}{fundamental_line_integrals}
If $\vv{F} = \nabla f$ on $D$ and $C$ is a piecewise smooth oriented curve in $D$ from $P$ to $Q$, then
\[
    \boxed{\int_C\vv{F}\cdot\dd{\vv{r}} = f(Q) - f(P).}
\]
Only the endpoint values of the potential matter.
\end{theorem}
For a parametrization $\vv{r}(t)$, $a \le t \le b$, the chain rule gives
\[
    \nabla f(\vv{r}(t))\cdot\vv{r}'(t) = \frac{\dd{}}{\dd{t}}f(\vv{r}(t)).
\]
Integrating and applying the ordinary Fundamental Theorem of Calculus proves the result.

\begin{example}[Reading a Potential Contour Diagram]
\begin{minipage}[c]{0.55\linewidth}
\vspace{1em}
In the diagram, $P$ lies on $f = 2$ and $Q$ on $f = 8$. Every path from $P$ to $Q$ gives
\[
    \int_C\vv{F}\cdot\dd{\vv{r}} = 8 - 2 = \boxed{6}.
\]
Read the endpoint contours, rather than counting crossings. Reversing the endpoints gives $-6$.
\end{minipage}\hfill
\begin{minipage}[c]{0.41\linewidth}
\centering
\begin{tikzpicture}[>=Stealth,font=\small,x=0.65cm,y=0.4cm]
    % Use schematic level sets to emphasize endpoint values rather than path shape.
    \foreach \y in {2,4,6,8}{
        \draw[green!45!black,thin] (-0.5,\y) -- (4,\y) node[right] {$f = \y$};
    }
    \draw[blue!70!black,thick,->] (0,2) .. controls (-1,6) and (1,9) .. (3,8);
    \draw[magenta!80!black,thick,->] (0,2) .. controls (3,1) and (4,5) .. (3,8);
    \fill (0,2) circle (1.8pt) node[below left] {$P$};
    \fill (3,8) circle (1.8pt) node[above right] {$Q$};
\end{tikzpicture}
\end{minipage}
\end{example}

\newpage
\subsection{Path Independence}
\begin{definition}{Path Independence}{path_independence}
A vector field is \textbf{path independent} on $D$ if its vector line integral has the same value along any two piecewise smooth paths in $D$ with the same starting and ending points.
\end{definition}

\begin{theorem}{Equivalent Characterizations}{conservative_equivalences}
For a continuous vector field on an open connected domain $D$, the following statements are equivalent:
\begin{enumerate}
    \item $\vv{F}$ is conservative on $D$. \vspace{0.5em}
    \item $\displaystyle\int_C\vv{F}\cdot\dd{\vv{r}}$ is independent of the path in $D$. \vspace{0.5em}
    \item $\displaystyle\oint_C\vv{F}\cdot\dd{\vv{r}} = 0$ for every closed piecewise smooth curve in $D$. \vspace{0.5em}
\end{enumerate}

Here, $\displaystyle\oint_C$ denotes the line integral around a closed curve $C$.
\end{theorem}

If $C_1$ and $C_2$ both run from $P$ to $Q$, then $C_1 + (-C_2)$ is closed. Thus zero circulation around every closed curve gives
\[
    0 = \int_{C_1}\vv{F}\cdot\dd{\vv{r}} - \int_{C_2}\vv{F}\cdot\dd{\vv{r}}.
\]
Conversely, path independence makes the integral around a closed path equal to the integral around the constant path, namely $0$.

\subsubsection{Constructing a Potential}

Suppose $\vv{F}$ is a continuous and path independent vector field on an open connected
domain $\mathcal{D}$, so that $\vv{F} = \nabla f$ for some function $f$. Fix a base point $P_0 \in \mathcal{D}$.

For each $P \in \mathcal{D}$, choose a piecewise smooth path $C_P$ from $P_0$ to $P$
and define
\[
    f(P) = \int_{C_P} \vv{F}\cdot\dd{\vv{r}}.
\]
Connectedness ensures that such a path exists. Path independence ensures
that every choice of $C_P$ gives the same value, so $f(P)$ is well defined.

\newpage
\noindent\begin{minipage}[c]{0.62\linewidth}
To compute $f_x(P)$, let $h \ne 0$ be small enough that the entire segment
from $P$ to $P + h\vv{i}$ lies in $\mathcal{D}$. Such a segment exists
because $\mathcal{D}$ is open.

To evaluate $f(P + h\vv{i})$, first follow
$C_P$ to $P$, then follow the short segment $S_h$ to $P + h\vv{i}$.
\end{minipage}\hfill
\begin{minipage}[c]{0.34\linewidth}
\centering\small
\begin{tikzpicture}[>=Stealth,font=\small]
\coordinate (A) at (0,0);
\coordinate (B) at (2.2,0.5);
\coordinate (H) at (3.5,0.5);
% The short segment extends the shared path.
\draw[blue!70!black,very thick,->] (A) to[out=65,in=155] node[above] {$C_P$} (B);
\draw[orange!85!black,very thick,->] (B) -- node[above] {$S_h$} (H);
\fill (A) circle (1.6pt) node[below] {$P_0$};
\fill (B) circle (1.6pt) node[below] {$P$};
\fill (H) circle (1.6pt) node[below] {$P + h\vv{i}$};
\end{tikzpicture}\par\smallskip
The blue contribution cancels. Shown for $h > 0$.
\end{minipage}\par

By additivity,
\[
    f(P + h\vv{i}) = f(P) + \int_{S_h}\vv{F}\cdot\dd{\vv{r}}.
\]
Parametrize the segment by $\vv{r}(t) = P + th\vv{i}$,
$0 \le t \le 1$. Its velocity is $h\vv{i}$, so the dot product
selects the $x$-component of $\vv{F}$ and multiplies it by $h$:
\[
    f(P + h\vv{i}) - f(P)
    = h\int_0^1 F_1(P + th\vv{i})\,\dd{t}.
\]
This parametrization works for both signs of $h$. Dividing by $h$ gives
\[
    \frac{f(P + h\vv{i}) - f(P)}{h}
    = \int_0^1 F_1(P + th\vv{i})\,\dd{t}.
\]
Every point $P + th\vv{i}$ is within distance $|h|$ of $P$. By continuity,
all the values of $F_1$ in this average approach $F_1(P)$ as $h \to 0$.
Therefore $f_x(P) = F_1(P)$. Taking short segments in the other
coordinate directions gives the remaining partial derivatives, so
$\nabla f = \vv{F}$.

\noindent\begin{minipage}[c]{0.62\linewidth}
\textbf{Changing the base point.}
If we start at $Q_0$ instead, we can reach $P$ by first going from $Q_0$
to $P_0$ and then following $C_P$. The new potential is
\[
    \widetilde f(P)
    = \int_{Q_0}^{P_0}\vv{F}\cdot\dd{\vv{r}} + f(P).
\]
The first term depends only on the two fixed base points. Thus changing
the base point adds the same constant to every value of $f$.
\end{minipage}\hfill
\begin{minipage}[c]{0.34\linewidth}
\centering\small
\begin{tikzpicture}[>=Stealth,font=\small]
\coordinate (Q) at (0,0);
\coordinate (A) at (1.4,0);
\coordinate (B) at (3.5,0.55);
% The initial path is fixed as P varies.
\draw[orange!85!black,very thick,->] (Q) -- node[above] {fixed path} (A);
\draw[blue!70!black,very thick,->] (A) to[out=65,in=155] node[above] {$C_P$} (B);
\fill (Q) circle (1.6pt) node[below] {$Q_0$};
\fill (A) circle (1.6pt) node[below] {$P_0$};
\fill (B) circle (1.6pt) node[below] {$P$};
\end{tikzpicture}\par\smallskip
The orange contribution is fixed.
\end{minipage}\par

\newpage
\subsubsection{Choosing an Easier Path}
\begin{problem}{A Complicated Component That Need Not Be Integrated}{easy_path}
The field $\vv{F} = \ang{x^{17}e^{x^2},7}$ is path independent. Let $C$ be the right half of $(x - 3)^2 + (y - 5)^2 = 4$, from $(3,3)$ to $(3,7)$. Evaluate $\int_C\vv{F}\cdot\dd{\vv{r}}$.
\end{problem}
\begin{sol}
Replace the semicircle by the vertical segment $\vv{r}(t) = \ang{3,t}$, $3 \le t \le 7$. Along this segment $\dd{x} = 0$, so
\[
    \int_C\vv{F}\cdot\dd{\vv{r}}
    = \int_3^7 7\,\dd{y} = \boxed{28}.
\]
There is no need to find an antiderivative of $x^{17}e^{x^2}$.
\end{sol}
\begin{center}
\begin{tikzpicture}[>=Stealth,font=\small,scale=0.65]
    % Both paths connect the same endpoints in the field's domain.
    \draw[blue!70!black,thick] (0,-2) arc (-90:90:2);
    \draw[->,blue!70!black,thick] (2,0) arc (0:25:2);
    \draw[->,orange!85!black,thick] (0,-2) -- (0,2);
    \fill (0,-2) circle (2pt) node[left] {$(3,3)$};
    \fill (0,2) circle (2pt) node[left] {$(3,7)$};
    \node[right] at (2,0) {$C$};
    \node[anchor=west] at (3,0) {Use the vertical shortcut.};
\end{tikzpicture}
\end{center}

In the three-path example from Section~16, $\vv{F} = \ang{y + 3x^2,x} = \nabla(xy + x^3)$. The endpoint difference from $(2,0)$ to $(-2,0)$ is $-8 - 8 = -16$, which accounts for both curved paths. The straight segment runs in reverse and gives $16$.

\begin{problem}{Does a Conservative Field Have Path-Independent Flux?}{flux_not_work}
If $\vv{F}$ is conservative and $C_1,C_2$ run from $A$ to $B$, must the flux across the two curves be the same?
\end{problem}
\begin{sol}
\textbf{False.} Conservativeness controls $\int_C\vv{F}\cdot\dd{\vv{r}}$, not $\int_C\vv{F}\cdot\vv{n}\,\dd{s}$.
For $\vv{F} = \ang{x,0} = \nabla(x^2/2)$, the right-normal flux is $\int_C x\,\dd{y}$. From $(0,0)$ to $(1,1)$, the diagonal $x = y = t$ gives $\int_0^1 t\,\dd{t} = 1/2$. Going first horizontally to $(1,0)$ and then vertically to $(1,1)$ gives $0 + 1 = 1$.
\end{sol}

\newpage
\subsection{Simply Connected Domains}
\begin{definition}{Simply Connected Domain}{simply_connected}
A domain $D$ is \textbf{simply connected} if it is path connected and every closed curve in $D$ can be continuously contracted to a point while staying in $D$.
\end{definition}
In the plane, a hole can obstruct this contraction. In three dimensions, removing a point and removing an entire line have different effects.
\begin{center}
\begin{tikzpicture}[>=Stealth,font=\small]
    % Contrast an annulus with the same band opened by a radial gap.
    \begin{scope}
        \fill[cyan!20,even odd rule] (0,0) circle (1.3) (0,0) circle (0.6);
        \draw[cyan!60!black] (0,0) circle (1.3) (0,0) circle (0.6);
        \draw[blue!70!black,thick] (0,0) circle (0.95);
        \draw[->,blue!70!black,thick] (0.95,0) arc (0:35:0.95);
        \node at (0,-1.8) {A: annulus --- no};
    \end{scope}
    \begin{scope}[shift={(5,0)}]
        \filldraw[fill=cyan!20,draw=cyan!60!black]
            (35:1.3) arc (35:325:1.3) -- (325:0.6) arc (325:35:0.6) -- cycle;
        \node at (0,-1.8) {B: open C-shaped band --- yes};
        \node[align=center] at (0,0) {gap};
    \end{scope}
\end{tikzpicture}
\end{center}
\begin{problem}{Which Domains Are Simply Connected?}{domain_classification}
Classify the two planar regions above and the following three-dimensional domains:
\[
    \text{C: }\mathbb{R}^3\setminus\{(0,0,0)\},\qquad
    \text{D: }\{(x,y,z):x^2 + y^2 > 0\}.
\]
\end{problem}
\begin{sol}
\textbf{B and C are simply connected; A and D are not.}
\begin{itemize}
    \item \textbf{A:} A loop winding around the annulus's hole cannot contract within the annulus.
    \item \textbf{B:} The gap opens the ring. Every loop can contract within the band.
    \item \textbf{C:} A loop can move through the third dimension and contract without meeting the missing point. Equivalently, the punctured space retracts onto the simply connected sphere $S^2$.
    \item \textbf{D:} This removes the entire $z$-axis. A loop winding around that axis cannot contract while avoiding it.
\end{itemize}
\end{sol}

\newpage
\subsection{Testing Whether a Field Is Conservative}
\begin{theorem}{Cross-Partials Test}{cross_partials_test}
Let $\vv{F}$ have continuous first partial derivatives on an open domain $D$.
\begin{itemize}
    \item If $\vv{F}$ is conservative, its cross partials agree:
    \[
        P_y = Q_x\quad\text{in two dimensions},\qquad
        \nabla\times\vv{F} = \vv{0}\quad\text{in three dimensions}.
    \]
    \item If $D$ is also simply connected, these conditions are sufficient for $\vv{F}$ to be conservative.
\end{itemize}
\end{theorem}
For $\vv{F} = \ang{P,Q,R}$, zero curl means all three equalities
\[
    P_y = Q_x,\qquad P_z = R_x,\qquad Q_z = R_y.
\]
On a domain with a hole, passing the derivative test alone does not settle conservativeness. Failing the test at even one point rules out a potential on the whole domain.

\begin{center}
\begin{tikzpicture}[>=Stealth,font=\small,
    box/.style={draw,rounded corners,align=center,inner sep=7pt}]
    % Separate unconditional equivalences from the converse requiring topology.
    \node[box,draw=blue!70!black] (a) at (0,0) {$\vv{F} = \nabla f$};
    \node[box] (b) at (4.7,0) {path independence};
    \node[box] (c) at (4.7,-2) {$\displaystyle\oint_C\vv{F}\cdot\dd{\vv{r}} = 0$\\for every closed $C$};
    \node[box,draw=red!70!black] (d) at (0,-2) {cross-partials\\condition};
    \draw[<->,thick] (a) -- (b);
    \draw[<->,thick] (b) -- (c);
    \draw[->,thick] (a) -- (d);
    \draw[->,dashed,thick] (d.west) -- ++(-1.2,0) |- (a.west);
    \node[align=center] at (2.1,-3.1) {Dashed converse: $D$ must be simply connected.};
\end{tikzpicture}
\end{center}
The equivalences use an open connected domain and a continuous field; the derivative statements additionally use continuous first partial derivatives.

For the rotational field from Section~15, the nonzero curl rules out a potential.

\newpage
\subsection{The Angle Field and the Importance of the Domain}
\begin{problem}{The Arctangent Warmup}{arctangent_domain}
Let $f(x,y) = \arctan(y/x)$. Find $\nabla f$ and examine the proposed integral $\oint_C\nabla f\cdot\dd{\vv{r}}$ around the counterclockwise unit circle.
\end{problem}
\begin{sol}
For $x \ne 0$, the chain rule gives
\[
    f_x = \frac{-y/x^2}{1 + y^2/x^2} = \frac{-y}{x^2 + y^2},\qquad
    f_y = \frac{1/x}{1 + y^2/x^2} = \frac{x}{x^2 + y^2}.
\]
However, $f$ and its gradient are undefined at the circle's points $(0,1)$ and $(0,-1)$. Thus the proposed integral of this gradient over the entire circle is \textbf{not defined in the usual sense}. The fundamental theorem cannot be applied there.

The formula for the gradient extends to a different field on the punctured plane:
\[
    \vv{G}(x,y) = \ang{\frac{-y}{x^2 + y^2},\frac{x}{x^2 + y^2}},\qquad
    D = \mathbb{R}^2\setminus\{(0,0)\}.
\]
For $\vv{r}(t) = \ang{\cos t,\sin t}$, $0 \le t \le 2\pi$,
\[
    \vv{G}(\vv{r}(t)) = \ang{-\sin t,\cos t} = \vv{r}'(t),\qquad
    \oint_C\vv{G}\cdot\dd{\vv{r}} = \int_0^{2\pi}1\,\dd{t} = \boxed{2\pi}.
\]
So $\vv{G}$ is not conservative on $D$. Nevertheless,
\[
    \pdv{G_1}{y} = \frac{y^2 - x^2}{(x^2 + y^2)^2}
    = \pdv{G_2}{x}
\]
throughout $D$. The missing origin prevents the curl-free test from implying conservativeness on this domain.
\end{sol}

\newpage
\subsubsection{A Local Angle Is Not a Global Potential}
\begin{center}
\begin{tikzpicture}[>=Stealth,font=\small,scale=1.4]
    % Show circulation on the punctured plane without assigning a global angle.
    \foreach \x in {-1.5,-0.75,0,0.75,1.5}{
        \foreach \y in {-1.5,-0.75,0,0.75,1.5}{
            \pgfmathsetmacro{\rsq}{(\x)^2 + (\y)^2}
            \ifdim\rsq pt>0pt
                \draw[->,gray!75] (\x,\y) -- ++({-0.3 * \y / \rsq},{0.3 * \x / \rsq});
            \fi
        }
    }
    \draw[->,black] (-2,0) -- (2,0) node[right] {$x$};
    \draw[->,black] (0,-2) -- (0,2) node[above] {$y$};
    \filldraw[fill=white,draw=red!75!black] (0,0) circle (0.06);
    \draw[blue!70!black,thick] (0,0) circle (1);
    \draw[->,blue!70!black,thick] (1,0) arc (0:45:1);
    \node[anchor=west] at (2.3,0.7) {$\displaystyle\oint_C\vv{G}\cdot\dd{\vv{r}} = 2\pi$};
    \node[anchor=west,align=left] at (2.3,-0.5) {One full turn changes\\a continuously tracked angle by $2\pi$.};
\end{tikzpicture}
\end{center}
Locally one can choose an angle $\theta$ with $\nabla\theta = \vv{G}$. On the whole punctured plane there is no single-valued differentiable angle potential: after one full turn the point is the same, but a continuously tracked angle has increased by $2\pi$.

For example, on the slit plane obtained by removing the nonpositive $x$-axis, a continuous branch $-\pi < \theta < \pi$ is a valid potential. A full circle is not contained in that domain. Similarly, $\arctan(y/x)$ is a potential separately on each half-plane $x > 0$ and $x < 0$.

\begin{tcolorbox}[colback=orange!5!white,colframe=orange!80!black,title=Check the Domain First]
A potential must be a single-valued function defined on a domain containing the \emph{entire} curve. Writing a local formula $\vv{G} = \nabla\theta$ does not prove that such a global potential exists.
\end{tcolorbox}

\newpage
\subsection{Finding a Potential Function}
For a planar field $\vv{F} = \ang{P,Q}$, first check the domain and $P_y = Q_x$. To construct a potential:
\begin{enumerate}
    \item Integrate $f_x = P$ with respect to $x$, treating $y$ as constant:
    \[
        f(x,y) = \int P(x,y)\,\dd{x} + g(y).
    \]
    \item Differentiate this expression with respect to $y$ and set $f_y = Q$. Solve for $g'(y)$, then integrate.
    \item Verify both components of $\nabla f$. In three dimensions, start with an unknown function $g(y,z)$ and match both remaining components.
\end{enumerate}

\begin{problem}{Work along a Spiral}{spiral_potential}
Let $\vv{F} = \ang{2xy + 3x^2,x^2 + \cos y}$ and
\[
    \vv{r}(t) = \ang{\sqrt{t}\cos(2\pi t),\sqrt{t}\sin(2\pi t)},\qquad 0 \le t \le 9.
\]
Evaluate $\int_C\vv{F}\cdot\dd{\vv{r}}$.
\end{problem}
\begin{sol}
On $\mathbb{R}^2$, $P_y = 2x = Q_x$, so the field is conservative. Integrate the first component:
\[
    f(x,y) = x^2y + x^3 + g(y),\qquad
    f_y = x^2 + g'(y) = x^2 + \cos y.
\]
Thus $g(y) = \sin y + K$, and a potential is
\[
    f(x,y) = x^2y + x^3 + \sin y.
\]
The endpoints are $\vv{r}(0) = (0,0)$ and $\vv{r}(9) = (3,0)$. Hence
\[
    \int_C\vv{F}\cdot\dd{\vv{r}} = f(3,0) - f(0,0) = \boxed{27}.
\]
\end{sol}

\newpage
\subsubsection{The Spiral: Geometry versus Computation}
\begin{center}
\begin{tikzpicture}[>=Stealth,font=\small,scale=0.7]
    % Use radius as parameter to keep the origin regular and retain all nine turns.
    \draw[->,black] (-3.5,0) -- (3.8,0) node[right] {$x$};
    \draw[->,black] (0,-3.5) -- (0,3.5) node[above] {$y$};
    \draw[blue!70!black,thin,domain=0:3,samples=1200,variable=\u]
        plot ({\u * cos(360 * \u^2)},{\u * sin(360 * \u^2)});
    \fill[red!75!black] (0,0) circle (2pt) node[above left] {$(0,0)$};
    \fill[red!75!black] (3,0) circle (2pt) node[below right] {$(3,0)$};
\end{tikzpicture}\par
\small Nine turns, but only two endpoint values are needed.
\end{center}
Although the original speed is unbounded at $t = 0$, substituting $u = \sqrt{t}$ gives a smooth regular parametrization:
\[
    \widetilde{\vv{r}}(u) = \ang{u\cos(2\pi u^2),u\sin(2\pi u^2)},\qquad 0 \le u \le 3.
\]
The fundamental theorem therefore applies to the geometric curve.

\subsection{When to Use an Endpoint Shortcut}
With a known potential, compute $f(Q) - f(P)$. With path independence already established, an easier path in the domain is another option. To justify either shortcut, verify a potential directly or check the cross partials together with the domain. For direct computation, use the work formula in the line integral reference in Section~16.
